%%=============================================================================
%% Samenvatting - VERPLICHT IN TWEE TALEN
%%=============================================================================
%%
%% Je rapport begint met twee samenvattingen, elk op een eigen bladzijde:
%%
%%   bladzijde 1: de Nederlandse samenvatting
%%   bladzijde 2: dezelfde samenvatting in het Engels
%%
%% Beide zijn verplicht, ook als je je rapport volledig in het Nederlands
%% schrijft. Ze staan bovenaan omdat ze in de praktijk het enige zijn dat een
%% druk jurylid met zekerheid leest.
%%
%% Richtlengte: 200 tot 250 woorden per taal, doorlopende tekst, geen lijstjes.
%%
%% Een goede samenvatting beantwoordt in die volgorde:
%%   1. Voor welk bedrijf en welk probleem?
%%   2. Wat heb je gemaakt?
%%   3. Met welke technologie, en waarom die?
%%   4. Wat is het resultaat - werkt het, en hoe weet je dat?
%%   5. Wat heeft het bedrijf eraan?
%%
%% Schrijf ze op het einde, niet in het begin. En schrijf ze zelf: dit is de
%% bladzijde waaraan een lezer merkt of je je eigen werk begrepen hebt.

%%---------- Nederlandse samenvatting -----------------------------------------

\IfLanguageName{english}{\selectlanguage{dutch}}{}

\chapter*{Samenvatting}
\label{ch:samenvatting-nl}

Vervang deze tekst door je Nederlandse samenvatting van ongeveer 200 woorden.

\IfLanguageName{english}{\selectlanguage{english}}{}

%%---------- Engelse samenvatting ---------------------------------------------

\IfLanguageName{dutch}{\selectlanguage{english}}{}

\chapter*{Abstract}
\label{ch:samenvatting-en}

Replace this text with the English version of your summary, about 200 words.
It is a translation of the Dutch summary, not a different text.

\IfLanguageName{dutch}{\selectlanguage{dutch}}{}
